% !TEX program = xelatex
\documentclass[11pt,a4paper]{article}

\usepackage[UTF8]{ctex}
\usepackage[margin=2.2cm]{geometry}
\usepackage{amsmath, amssymb, amsthm}
\usepackage{listings}
\usepackage{xcolor}
\usepackage{tikz}
\usepackage{booktabs}
\usepackage{multirow}
\usepackage{hyperref}
\usetikzlibrary{positioning, arrows.meta, shapes.geometric, fit, calc}

\hypersetup{
  colorlinks=true,
  linkcolor=blue!60!black,
  urlcolor=teal!70!black,
}

\lstdefinelanguage{Rust}{
  morekeywords={fn,let,mut,pub,struct,impl,use,mod,self,Self,Option,Some,None,return,if,else,for,while,loop,match,break,continue,as,where,trait,type,enum,move,ref,const,static,unsafe,extern,crate,super,dyn,async,await,box},
  sensitive=true,
  morecomment=[l]{//},
  morecomment=[s]{/*}{*/},
  morestring=[b]",
}

\lstset{
  language=Rust,
  basicstyle=\ttfamily\footnotesize,
  keywordstyle=\color{blue!60!black}\bfseries,
  commentstyle=\color{green!40!black}\itshape,
  stringstyle=\color{red!60!black},
  numbers=left,
  numberstyle=\tiny\color{gray},
  numbersep=8pt,
  frame=single,
  framerule=0.4pt,
  rulecolor=\color{gray!40},
  breaklines=true,
  breakatwhitespace=true,
  showstringspaces=false,
  tabsize=4,
  columns=fullflexible,
  keepspaces=true,
}

\title{\textbf{halfedge\,: 底层存储模块设计文档} \\ \large \texttt{src/storage.rs}}
\author{halfedge 库}
\date{2026-07-04}

\begin{document}
\maketitle
\tableofcontents

\section{定位与边界}

\texttt{MeshStorage} 是整个半边网格库的\textbf{底座}，只承担两类职责：
\begin{enumerate}
  \item 三类元素（顶点 / 半边 / 面）的\textbf{存储与生命周期管理}；
  \item 句柄\textbf{有效性判定}与安全访问。
\end{enumerate}

它\textbf{不}负责维护任何拓扑一致性约束——\texttt{twin}、\texttt{next}、\texttt{prev}、
\texttt{face} 等关系的正确性由上层构建器保证。这样切分使得底座极其稳定、可独立测试，
后续可以叠加不同策略的构建器而不影响存储层。

\section{数据结构}

\subsection{三类元素}

\begin{lstlisting}
pub struct Vertex {
    pub position: [f64; 3],
    pub halfedge: Option<HalfEdgeId>, // 从该顶点出发的任一 outgoing 半边
}

pub struct HalfEdge {
    pub vertex: VertexId,              // 半边指向的目的顶点 (tip)
    pub twin:  Option<HalfEdgeId>,
    pub next:  Option<HalfEdgeId>,
    pub prev:  Option<HalfEdgeId>,
    pub face:  Option<FaceId>,
}

pub struct Face {
    pub halfedge: Option<HalfEdgeId>,  // 边界环入口
}
\end{lstlisting}

三类元素均提供 \texttt{new()} 构造函数。\texttt{Vertex} 与 \texttt{Face}
额外实现 \texttt{Default} trait：
\begin{itemize}
  \item \texttt{Vertex::default()} $\equiv$ \texttt{Vertex::new([0.0; 3])}（原点 + 无半边）；
  \item \texttt{Face::default()} $\equiv$ \texttt{Face::new()}（无半边）。
\end{itemize}
这使得二者可在需要 \texttt{Default} bound 的泛型上下文中使用
（如 \texttt{SlotMap::insert\_with}、容器预分配初始化等）。
\texttt{HalfEdge} 因必须指定 \texttt{vertex} 字段而无 \texttt{Default}。

\subsection{半边四向邻接示意}

半边 $h$ 通过四个指针把网格「编织」起来，下图展示一个三角面内的连接关系：

\begin{center}
\begin{tikzpicture}[
  vert/.style={circle, draw=blue!70!black, fill=blue!10, minimum size=7mm, inner sep=0pt, font=\small},
  he/.style={-Stealth, thick, draw=red!70!black},
  he lab/.style={font=\footnotesize, sloped, above},
  face/.style={regular polygon, regular polygon sides=3, draw=gray!60, fill=gray!8, minimum size=3.6cm, rotate=0}
]
  \node[face] (F) at (0,0) {};
  \node[vert] (v0) at (F.corner 1) {$v_0$};
  \node[vert] (v1) at (F.corner 2) {$v_1$};
  \node[vert] (v2) at (F.corner 3) {$v_2$};

  \draw[he] (v0) -- node[he lab] {$h_0$} (v1);
  \draw[he] (v1) -- node[he lab] {$h_1$} (v2);
  \draw[he] (v2) -- node[he lab] {$h_2$} (v0);

  \node[font=\small, text=gray!50!black] at (0,-0.05) {$F$};

  \node[below=2.2cm of v1, font=\footnotesize, align=center] (legend)
    {$h_i.\mathrm{next}$: $h_0{\to}h_1{\to}h_2{\to}h_0$\quad
     $h_i.\mathrm{prev}$: 反向环\quad
     $h_i.\mathrm{face}=F$\quad
     $h_i.\mathrm{vertex}=\text{tip}$};
\end{tikzpicture}
\end{center}

对应的形式化关系（同面边界环）：
\[
\forall i:\quad h_i.\mathrm{face}=F,\quad
h_i.\mathrm{next}=h_{(i+1)\bmod 3},\quad
h_i.\mathrm{prev}=h_{(i-1)\bmod 3},\quad
h_i.\mathrm{vertex}=v_{(i+1)\bmod 3}.
\]

\texttt{twin} 跨面连接：若两条半边共享两端点但方向相反，则互为 twin。
边界半边的 \texttt{twin} 与 \texttt{face} 通常为 \texttt{None}。

\subsection{容器}

\begin{lstlisting}
pub struct MeshStorage {
    vertices:  SlotMap<VertexId,   Vertex>,
    halfedges: SlotMap<HalfEdgeId, HalfEdge>,
    faces:     SlotMap<FaceId,     Face>,
    // SOA 位置缓存（2026-07 新增，提升缓存命中率）
    positions:  Vec<[f64; 3]>,
    pos_index:  SecondaryMap<VertexId, u32>,
}
\end{lstlisting}

三类元素分别放在独立的 \texttt{SlotMap} 中，互不干扰。三组 \texttt{Id} 在编译期
已被 \texttt{new\_key\_type!} 区分（详见 \texttt{ids.tex}），即使误把
\texttt{HalfEdgeId} 传给 \texttt{get\_vertex} 也会编译失败。

\subsection{SOA 位置缓存}

\texttt{Vertex} 结构体内含 \texttt{position}（24 B）与 \texttt{halfedge}（8 B）
两个字段，加上 SlotMap 槽位元数据（4 B 版本号），单槽步长约 36 B。
几何热路径（\texttt{mesh\_volume}、\texttt{mesh\_aabb}、\texttt{ray\_mesh\_intersects}
等）仅访问 \texttt{position}，24 B 有效数据在 36 B 槽位中仅占 67\%，造成缓存浪费。

为提升缓存命中率，\texttt{MeshStorage} 内部并行维护一份稠密的
\texttt{positions: Vec<[f64; 3]>} 缓存，步长 24 B，利用率 100\%。
\texttt{pos\_index: SecondaryMap<VertexId, u32>} 提供
\texttt{VertexId $\to$ u32} 的稠密索引映射（SecondaryMap 由 Vec 支持，$O(1)$ 查找）。

\begin{center}
\begin{tikzpicture}[
  node distance=0.6cm and 1.0cm,
  blk/.style={draw, rounded corners=2pt, minimum width=3.4cm, minimum height=0.9cm, align=center, font=\small},
  arr/.style={-Stealth, thick}
]
  \node[blk, fill=blue!10] (slotmap) {SlotMap<VertexId, Vertex>\\36 B/槽};
  \node[blk, fill=green!10, right=of slotmap] (vec) {Vec<[f64;3]>\\24 B/槽};
  \node[blk, fill=orange!10, below=of vec] (idx) {SecondaryMap<VertexId, u32>\\8 B/槽};
  \draw[arr] (slotmap.east) -- node[above, font=\scriptsize] {同步} (vec.west);
  \draw[arr, dashed] (slotmap.south) |- (idx.west);
  \draw[arr, dashed] (idx.north) -- node[right, font=\scriptsize] {索引} (vec.south);
\end{tikzpicture}
\end{center}

\textbf{同步策略}：
\begin{sloppypar}
\begin{itemize}
  \item \texttt{add\_vertex}：同时推入 \texttt{positions} 与 \texttt{pos\_index}；
  \item \texttt{remove\_vertex}：swap-remove 位置并更新填补元素的索引；
  \item \texttt{set\_position}：同时更新主存与缓存（推荐的位置修改入口）；
  \item \texttt{get\_vertex\_mut}：若直接修改 \texttt{position}，缓存会暂时失效，需调用 \texttt{sync\_position} 或 \texttt{rebuild\_position\_cache} 修复。
\end{itemize}
\end{sloppypar}

\textbf{公开 API}：
\begin{center}
\renewcommand{\arraystretch}{1.15}
\begin{tabular}{@{}lll@{}}
  \toprule
  \textbf{方法} & \textbf{返回} & \textbf{用途} \\
  \midrule
  \texttt{get\_position(id)}        & \texttt{Option<[f64;3]>} & 读 SOA 缓存取单点位置 \\
  \texttt{positions\_dense()}        & \texttt{\&[[f64;3]]}    & 批量遍历稠密切片 \\
  \texttt{position\_index(id)}        & \texttt{Option<u32>}    & VertexId $\to$ 稠密索引 \\
  \texttt{set\_position(id, pos)}    & \texttt{Option<()>}     & 同步更新主存与缓存 \\
  \texttt{sync\_position(id)}        & \texttt{()}             & 单点同步（修复 get\_vertex\_mut 的副作用） \\
  \texttt{rebuild\_position\_cache()} & \texttt{()}             & 全量重建缓存（反序列化后） \\
  \bottomrule
\end{tabular}
\end{center}

\textbf{收益预期}：顺序遍历顶点位置时（如 \texttt{mesh\_aabb}、\texttt{mesh\_centroid}），
缓存利用率从 67\% 提升至 100\%，缓存行覆盖顶点数从 $\lfloor 64/36 \rfloor = 1$ 提升至
$\lfloor 64/24 \rfloor = 2$。随机访问（如 \texttt{mesh\_volume} 按面取顶点）的改善
取决于顶点 ID 空间局部性，通常 10\%--30\%。

\section{核心 API 流程}

\subsection{新增}

\begin{center}
\begin{tikzpicture}[
  node distance=0.7cm,
  stage/.style={draw, rounded corners=2pt, minimum width=3.2cm, minimum height=0.7cm, align=center, font=\small},
  op/.style={-Stealth, thick}
]
  \node[stage, fill=green!12] (s1) {调用 \texttt{add\_vertex(v)}};
  \node[stage, fill=yellow!20, right=of s1] (s2) {查找空闲槽位 $i$};
  \node[stage, fill=yellow!20, right=of s2] (s3) {写入数据\\版本号 $v_i$};
  \node[stage, fill=blue!12, right=of s3] (s4) {返回 $(i, v_i)$};
  \draw[op] (s1) -- (s2);
  \draw[op] (s2) -- (s3);
  \draw[op] (s3) -- (s4);
\end{tikzpicture}
\end{center}

返回值即为外部使用的强类型 ID。三条 \texttt{add\_*} 接口复杂度均为 $O(1)$ 摊还。

\subsection{删除（自动失效）}

\begin{center}
\begin{tikzpicture}[
  node distance=0.8cm,
  stage/.style={draw, rounded corners=2pt, minimum width=3.6cm, minimum height=0.7cm, align=center, font=\small},
  op/.style={-Stealth, thick},
  dec/.style={diamond, draw, aspect=2, fill=orange!15, inner sep=1pt, font=\small, align=center}
]
  \node[stage, fill=green!12] (a) {调用 \texttt{remove\_vertex(id)}};
  \node[dec, below=of a] (b) {句柄有效？};
  \node[stage, fill=red!15, below left=of b, xshift=-1.0cm] (c) {清空槽位 $i$\\$v_i\mathrel{+}{=}1$};
  \node[stage, fill=gray!10, below right=of b, xshift=1.0cm] (d) {返回 \texttt{None}};
  \node[stage, fill=blue!12, below=of c] (e) {返回 \texttt{Some(旧值)}};

  \draw[op] (a) -- (b);
  \draw[op] (b) -- node[left, font=\footnotesize] {是} (c);
  \draw[op] (b) -- node[right, font=\footnotesize] {否} (d);
  \draw[op] (c) -- (e);
\end{tikzpicture}
\end{center}

\textbf{关键点：}版本号 $+1$ 是「自动标记失效」的核心。原句柄 $(i, v_i{-}1)$ 在后续
\texttt{get} 中因版本不匹配返回 \texttt{None}；即使槽位 $i$ 被 \texttt{add} 复用，
新元素的版本是 $v_i$（或更高），旧句柄仍然无效，从而\textbf{杜绝 ABA 问题}。

\subsection{访问（安全 \texttt{get}/\texttt{get\_mut}）}

所有访问方法返回 \texttt{Option}：
\begin{align*}
\texttt{get\_vertex(id)}      &: \texttt{VertexId} \to \texttt{Option<\&Vertex>} \\
\texttt{get\_vertex\_mut(id)} &: \texttt{VertexId} \to \texttt{Option<\&mut Vertex>}
\end{align*}

调用方\textbf{绝不会因为传入了已删除的 ID 而 panic}，只会得到 \texttt{None}。
配合 \texttt{contains\_*} 系列，可在不确定时先做有效性预检。

\section{API 一览}

\begin{center}
\renewcommand{\arraystretch}{1.18}
\begin{tabular}{@{}lll@{}}
  \toprule
  \textbf{类别} & \textbf{方法} & \textbf{返回值} \\
  \midrule
  构造 & \texttt{new()} & \texttt{MeshStorage} \\
  \midrule
  \multirow{3}{*}{增}
    & \texttt{add\_vertex(v)}    & \texttt{VertexId}   \\
    & \texttt{add\_halfedge(h)}  & \texttt{HalfEdgeId} \\
    & \texttt{add\_face(f)}      & \texttt{FaceId}     \\
  \midrule
  \multirow{3}{*}{删}
    & \texttt{remove\_vertex(id)}    & \texttt{Option<Vertex>}   \\
    & \texttt{remove\_halfedge(id)}  & \texttt{Option<HalfEdge>}  \\
    & \texttt{remove\_face(id)}      & \texttt{Option<Face>}      \\
  \midrule
  \multirow{3}{*}{查（不可变）}
    & \texttt{get\_vertex(id)}       & \texttt{Option<\&Vertex>}    \\
    & \texttt{get\_halfedge(id)}     & \texttt{Option<\&HalfEdge>}  \\
    & \texttt{get\_face(id)}         & \texttt{Option<\&Face>}      \\
  \midrule
  \multirow{3}{*}{查（可变）}
    & \texttt{get\_vertex\_mut(id)}  & \texttt{Option<\&mut Vertex>}   \\
    & \texttt{get\_halfedge\_mut(id)}& \texttt{Option<\&mut HalfEdge>} \\
    & \texttt{get\_face\_mut(id)}    & \texttt{Option<\&mut Face>}     \\
  \midrule
  \multirow{3}{*}{有效性}
    & \texttt{contains\_vertex(id)}    & \texttt{bool} \\
    & \texttt{contains\_halfedge(id)}  & \texttt{bool} \\
    & \texttt{contains\_face(id)}      & \texttt{bool} \\
  \midrule
  \multirow{3}{*}{统计}
    & \texttt{vertex\_count()}   & \texttt{usize} \\
    & \texttt{halfedge\_count()} & \texttt{usize} \\
    & \texttt{face\_count()}     & \texttt{usize} \\
  \midrule
  \multirow{2}{*}{容量管理}
    & \texttt{clear()}    & \texttt{()} \\
    & \texttt{reserve(v\_cap, he\_cap, f\_cap)} & \texttt{()} \\
  \bottomrule
\end{tabular}
\end{center}

\subsection{容量管理方法}

\begin{itemize}
  \item \texttt{clear(\&mut self)}：清空所有顶点、半边、面。等效于
        \texttt{*self = MeshStorage::new()}。\texttt{SlotMap::clear()} 是 $O(n)$
        时间但会释放内部内存，适合频繁重建网格的场景。
  \item \texttt{reserve(\&mut self, vertex\_cap, halfedge\_cap, face\_cap)}：
        为三类 \texttt{SlotMap} 预分配容量，减少批量构建时的 rehash / 重分配次数。
        三参数均为下限提示，实际分配可能更多。\texttt{io::build\_mesh\_*}
        构建器在入口处自动调用此方法。
\end{itemize}

\section{复杂度}

\begin{center}
\renewcommand{\arraystretch}{1.18}
\begin{tabular}{@{}ll@{}}
  \toprule
  \textbf{操作} & \textbf{复杂度} \\
  \midrule
  \texttt{add\_*}        & $O(1)$ 摊还（可能扩容） \\
  \texttt{remove\_*}     & $O(1)$ \\
  \texttt{get\_*} / \texttt{get\_*\_mut} & $O(1)$ \\
  \texttt{contains\_*}   & $O(1)$ \\
  \texttt{*\_count()}    & $O(1)$ \\
  \texttt{clear()}       & $O(n)$（释放内存） \\
  \texttt{reserve(...)}  & $O(n)$（可能重分配） \\
  \bottomrule
\end{tabular}
\end{center}

\section{单元测试覆盖}

\texttt{src/storage.rs} 末尾的 \texttt{\#[cfg(test)] mod tests} 共 54 个测试，
按需求点对应如下：

\begin{center}
\renewcommand{\arraystretch}{1.18}
\begin{tabular}{@{}ll@{}}
  \toprule
  \textbf{测试名} & \textbf{覆盖点} \\
  \midrule
  \multicolumn{2}{l}{\textit{基础 CRUD（3 个）}} \\
  \texttt{add\_and\_get\_vertex}     & 新增 + \texttt{get} + \texttt{contains} + 计数 \\
  \texttt{add\_and\_get\_halfedge}   & 半边新增 + 关联顶点 outgoing \\
  \texttt{add\_and\_get\_face}       & 面新增 \\
  \midrule
  \multicolumn{2}{l}{\textit{删除与 ABA 安全（4 个）}} \\
  \texttt{remove\_vertex\_invalidates\_id}    & 删除后 ID 失效 + 重复删除返回 \texttt{None} \\
  \texttt{remove\_halfedge\_invalidates\_id}  & 半边删除失效 \\
  \texttt{remove\_face\_invalidates\_id}      & 面删除失效 \\
  \texttt{slot\_reuse\_does\_not\_resurrect\_old\_id} & ABA 安全：槽位复用不复活旧句柄 \\
  \midrule
  \multicolumn{2}{l}{\textit{可变访问与无效句柄（2 个）}} \\
  \texttt{get\_mut\_allows\_in\_place\_update} & 可变访问就地修改拓扑字段 \\
  \texttt{accessing\_invalid\_id\_never\_panics} & 已删除 ID 访问绝不 panic \\
  \midrule
  \multicolumn{2}{l}{\textit{容量管理（5 个）}} \\
  \texttt{clear\_empties\_all\_three\_slotmaps} & \texttt{clear} 清空三类元素 + 迭代器为空 \\
  \texttt{clear\_equivalent\_to\_new} & \texttt{clear} 后与 \texttt{new()} 等价 \\
  \texttt{clear\_allows\_reuse\_after} & \texttt{clear} 后仍可继续添加元素 \\
  \texttt{reserve\_does\_not\_change\_counts} & \texttt{reserve} 仅预分配，不改变计数 \\
  \texttt{reserve\_zero\_is\_noop}    & \texttt{reserve(0,0,0)} 不 panic \\
  \midrule
  \multicolumn{2}{l}{\textit{迭代与空判断（5 个）}} \\
  \texttt{is\_empty\_on\_new\_mesh}   & 新建网格 \texttt{is\_empty()} 返回 true \\
  \texttt{is\_empty\_after\_clear}    & \texttt{clear} 后 \texttt{is\_empty()} 返回 true \\
  \texttt{vertices\_iter\_yields\_all\_data} & \texttt{vertices()} 迭代所有顶点数据 \\
  \texttt{vertices\_mut\_allows\_modification} & \texttt{vertices\_mut()} 允许就地修改 \\
  \texttt{empty\_iterators\_yield\_nothing} & 空网格所有数据迭代器不产出 \\
  \midrule
  \multicolumn{2}{l}{\textit{SOA 位置缓存（9 个，2026-07 新增）}} \\
  \texttt{soa\_cache\_get\_position\_matches\_vertex} & \texttt{get\_position} 与主存一致 \\
  \texttt{soa\_cache\_positions\_dense\_length\_matches} & 稠密切片长度与顶点数一致 \\
  \texttt{soa\_cache\_position\_index\_round\_trip} & VertexId $\to$ 索引 $\to$ 位置往返 \\
  \texttt{soa\_cache\_set\_position\_syncs\_both} & \texttt{set\_position} 同步主存与缓存 \\
  \texttt{soa\_cache\_remove\_preserves\_dense\_layout} & 删除后稠密布局保持正确 \\
  \texttt{soa\_cache\_sync\_position\_after\_get\_vertex\_mut} & \texttt{sync\_position} 修复缓存失效 \\
  \texttt{soa\_cache\_rebuild\_from\_master} & \texttt{rebuild\_position\_cache} 全量重建 \\
  \texttt{soa\_cache\_clear\_empties\_cache} & \texttt{clear} 同步清空缓存 \\
  \texttt{soa\_cache\_empty\_mesh\_no\_panic} & 空网格访问缓存不 panic \\
  \texttt{soa\_cache\_icosphere\_consistency} & icosphere 全顶点缓存一致性 \\
  \midrule
  \multicolumn{2}{l}{\textit{其他（含 Serde、拓扑诊断等，26 个）}} \\
  \multicolumn{2}{l}{含 \texttt{slot\_reuse\_after\_clear\_works}、\texttt{large\_icosphere\_counts\_correct}、} \\
  \multicolumn{2}{l}{\texttt{twin\_relationship\_is\_symmetric}、\texttt{serde\_roundtrip\_preserves\_topology} 等。} \\
  \bottomrule
\end{tabular}
\end{center}

全库 \texttt{cargo test} 当前共 669 个单元测试 + 24 个集成测试，全部通过。

\section{后续扩展}

存储层刻意保持极简，后续模块基于此构建了完整能力：
\begin{itemize}
  \item \textbf{拓扑操作}（\texttt{topology\_ops.rs}）：实现 \texttt{split\_edge}、
        \texttt{flip\_edge}、\texttt{collapse\_edge}、\texttt{extrude\_face}、
        \texttt{add\_triangle} 等高级操作，负责维护 \texttt{twin/next/prev} 一致性；
  \item \textbf{遍历器}（\texttt{traversal.rs}）：基于 \texttt{next/prev} 环遍历面边界，
        基于 \texttt{twin} 跨面遍历邻接面，提供 eager/lazy 两类迭代器与无向边迭代器；
  \item \textbf{序列化}（\texttt{io.rs}）：\texttt{OBJ/PLY} 加载与保存，支持 n-gon；
  \item \textbf{并行}：\texttt{SlotMap} 内部槽位互相独立，方便后续做并行构建。
\end{itemize}

\end{document}
